% Generated draft: requires mathematical and source review.
\documentclass[11pt]{article}
\usepackage[margin=25mm]{geometry}
\usepackage{fontspec}
\setmainfont{Latin Modern Roman}
\usepackage{amsmath,amssymb,mathtools,unicode-math}
\setmathfont{Latin Modern Math}
\usepackage{enumitem,needspace,xurl,hyperref}
\setlength{\parindent}{0pt}
\setlength{\parskip}{5pt}
\setlength{\emergencystretch}{3em}
\newcommand{\N}{\mathbb{N}}
\newcommand{\Z}{\mathbb{Z}}
\newcommand{\Q}{\mathbb{Q}}
\newcommand{\R}{\mathbb{R}}
\newcommand{\C}{\mathbb{C}}
\newcommand{\F}{\mathbb{F}}
\newcommand{\A}{\mathbb{A}}
\newcommand{\PP}{\mathbb{P}}
\newcommand{\E}{\mathbb{E}}
\newcommand{\eps}{\varepsilon}
\newcommand{\dd}{\,\mathrm{d}}
\begin{document}
\section*{UnsolvedMath ID: 20002314}
\textbf{Source problem code:} catalogue entry
\subsection{Definitions and mathematical statement}
Write \(I(T,\lambda)\) for the number of isomorphism types of \(T\)-models of cardinality \(\lambda\). In the intended countable-theory context, the source asks whether
\[
 \forall\lambda\geq\aleph_1\quad I(T,\lambda)\leq\lambda
\]
forces every such model to have a \(\lambda\)-computable presentation, allowing the extra hypothesis \(V=L\) if helpful. It says such theories are \(\omega\)-stable and non-multidimensional and records a positive \(\aleph_1\)-categorical case. Those surrounding theorem and computability conventions are not reproduced fully.

The aid gives a noncategorical special class. Fix \(m\geq2\), and let \(T_m\) be the complete theory of \(m\) named infinite unary predicates partitioning the universe, with no further structure. A model is classified by its finite dimension vector
\[
 (\kappa_0,\ldots,\kappa_{m-1}),\qquad
 \kappa_i\text{ infinite cardinals},\qquad\max_i\kappa_i=\lambda.
\]
For every uncountable \(\lambda\), the aid reports \(I(T_m,\lambda)\leq\lambda\). At \(\lambda=\aleph_1\), each coordinate is either \(\aleph_0\) or \(\aleph_1\), with at least one large coordinate, so
\[
 I(T_m,\aleph_1)=2^m-1.
\]
This is genuinely noncategorical for \(m\geq2\).

For every regular uncountable \(\lambda\), its computable presentation has domain the ordinals below \(\lambda\). Use disjoint columns
\[
 C_i=\{\omega\cdot\beta+i:\beta<\lambda\},\qquad0\leq i<m.
\]
Choose a designated \(j\) with \(\kappa_j=\lambda\). For \(i\ne j\), use all of \(C_i\) if \(\kappa_i=\lambda\), and only its terms with \(\beta<\kappa_i\) if \(\kappa_i<\lambda\). Assign every leftover ordinal to predicate \(j\). That predicate contains \(C_j\), so still has cardinality \(\lambda\); the construction defines a genuine partition and does not leave unused ordinals outside the model.

The aid reports the predicates uniformly \(\Delta_1(L_\lambda)\) from a finite ordinal index coding the small coordinates, the designated component and flags for \(\lambda\)-sized coordinates. This is an explicitly parameter/index-sensitive meaning of \(\lambda\)-computable; the stage ordinal \(\lambda\) is not used as an element of \(L_\lambda\). It is not a claim of parameter-free computation from \(m\) alone.

Thus the model-count hypothesis holds at all uncountable sizes in this class, and presentations are supplied at every regular uncountable size under the specified convention. No corresponding singular-cardinal presentation theorem, or complete answer for arbitrary \(T\), follows from the construction. The original all-\(\lambda\) scope and its undefined computational convention remain qualifications.
\subsection{Short English statement}
The aid supplies a noncategorical family satisfying the source's small-spectrum hypothesis: finitely many named infinite pure pieces. Every regular uncountable-sized model has an explicit ordinal-column presentation, with the finite dimension vector encoded in its allowed index. This is a special case and does not settle all theories, singular cardinals or parameter-free conventions.
\subsection{Sources}
\begin{itemize}
\item Local source record: catalogue entry, catalogue entry; canonical \texttt{aim-logic-notes.json}, \texttt{notes[89]}, canonical source item 1. Title: ``Computable copies for a finite dimension-vector spectrum.''
\item AIM list \emph{Computable stability theory}, Section 1.6, ``\(\lambda\)-many models of each cardinality \(\lambda\geq\aleph_1\) (Greenberg),'' supplied reference \url{https://aimath.org/pastworkshops/computestabproblems.pdf}. The extracted question number is split as 1 and 2; no unsupported printed-number repair is made here. This draft does not access a PDF.
\item Original contextual claims: \(\omega\)-stability, non-multidimensionality and the \(\aleph_1\)-categorical case. Attempt 1 supplies the finite-predicate family, spectrum and \(\Delta_1(L_\lambda)\) column construction, without full external bibliographic references.
\item Registry dated 14 August 2026: ``partially\_solved,'' ``PARTIAL-PROGRESS,'' statement status ``exact''; attempt partial, research-run status unknown, rigor ``proved,'' novelty confidence low. The machine-generated, non-peer-reviewed aid requires independent verification before citation.
\end{itemize}
\end{document}
